\defmodule {NORM}

This module provides tools to define normalization constants
for the shortest vectors.

\bigskip\hrule
\code

DEFINITION MODULE NORM;

IMPORT BASISLR;
FROM CONFIG IMPORT BDim, MDim;
FROM MULT   IMPORT MScal;

TYPE
   CriterionType = (Spectral, SpectralP, SpectralL1, Beyer, BeyerSpectral,
                    Palpha);
\endcode
 \tabb  Types of merit criteria for lattices.
   {\tt Spectral} means that the criterion is based on the shortest
   vector in the {\em dual\/} lattice.
   In this case, one must specify which approximation is used for $\gamma_t$
   (see {\tt GammaType} below).
   {\tt SpectralP} is similar to {\tt Spectral}, except that
    it is based on the shortest vector in the {\em primal\/} lattice.
   {\tt SpectralL1} is similar to {\tt Spectral}, except that
    the length of the (dual) vectors are measured with the $L_1$ norm,
    minus 1.  The length of the shortest dual vector is then an upper
    bound on the minimal number of hyperplanes that cover all the points
    (\cite{rDIE75a} and \cite{rKNU97a}, Exercises 3.3.4-15 and 16).
   {\tt Beyer} means that the criterion is the Beyer quotient $q_t$.
   {\tt BeyerSpectral} means that both {\tt Beyer} and {\tt Spectral}
   apply simultaneously.
   {\tt Palpha} means that the criterion is $P_\alpha$.
   In this case, one must specify a value of $\alpha$
   and weights may also be used (see the module {\tt PALPHA}).   
 \endtabb
\code

   NormType = (L1Norm, L2Norm, SupNorm);
\endcode
 \tabb  Type of norm used to measure the vector lengths.
  For $V = (v_1,\dots,v_t)$, the $L_1$ norm is 
  $\Vert V\Vert = |v_1|+\cdots+|v_t|$, the $L_2$ (or Euclidean) norm is
  $\Vert V\Vert = (v_1^2+\cdots+v_t^2)^{1/2}$, and the sup norm is
  $\Vert V\Vert = \max(|v_1|,\dots,|v_t|)$.
 \endtabb
\code

   GammaType = (BestLat, Laminated, Minkowski, Rogers, MinkL1);
\endcode
 \tabb  Types of bound on $\gamma_t$ which can be used to normalize 
   the length of the shortest vector.
   The normalized value is the  
   figure of merit used as a selection  criterion.
   {\tt BestLat, Laminated, Minkowski, Rogers} correspond
   respectively to the bounds $\gamma_t^{\rmL}$, $\gamma_t^{\rmB}$, 
   $\gamma_t^{\rmZ}$, and $\gamma_t^{\rmR}$, for the Euclidean norm.
   {\tt MinkL1} must be used for the $\cL_1$ norm and corresponds to 
   $\gamma_t^{\rmM} = (t!)^{1/t}$.
 \endtabb
\code

   Normaliz = POINTER TO InfoNormaliz;
   InfoNormaliz = RECORD
      Name  : ARRAY [0..31] OF CHAR;
      Norm  : NormType;
      GType : GammaType;
      Beta  : LONGREAL;
      Maxt  : BDim;
      Gamma : BASISLR.BVect;
      Cst   : BASISLR.BVect;
      END;
\endcode
  \tabb Data structure used to store a vector of values 
   that are used to normalize the lengths
   $\ell_t$ of the shortest vectors.  
   The array {\tt Gamma} contains values of $\tilde\gamma_t$ that are  
   approximations of $\gamma_t$ which appears in the bounds given
   in Section~1.
   The normalized figure of merit will be $\ell_t/\bar\ell_t$, where
  \eq             
    \bar\ell_t = \cases {
        \tilde\gamma_t m \beta^t       & if $t \le k$;\cr
        \tilde\gamma_t m^{k/t} \beta^t & if $t > k$,\rule{0pt}{14pt}\cr }
                                                    \eqlabel{eq:barellt}
  \endeq
   where $\beta$ is an additional multiplicative factor 
   stored in the variable {\tt Beta} (often equal to 1).
   The vector {\tt Cst} contains the values of $\bar\ell_t^2$ 
   in case of the Euclidean norm, and $\bar\ell_t$ otherwise.

   The structure must be created by the procedure {\tt CreateNormaliz},
   initialized by {\tt InitNormaliz}, and deleted by
   {\tt DeleteNormaliz}.
   The field {\tt Name} is a character string that describes the type of 
   $\tilde\gamma_t$.
   {\tt Norm} is the type of norm that is used to measure
   the vector lengths, {\tt GType} indicates the type of normalization 
   constants $\tilde\gamma_t$, 
   {\tt Beta} is the weighting factor used in (\ref{eq:barellt}),
   {\tt Maxt} is the maximum index for which the arrays have been
   initialized, i.e., the largest value of $t$ that can be considered.
%  Old: With {\tt Critere2 = MinkL1}, {\tt Gamma[t]} contains $t!$.
  \endtabb
\code

PROCEDURE CreateNormaliz (VAR G : Normaliz);
\endcode
 \tab  Creates the structure {\tt G} without initializing anything in it.
 \endtab
\code

PROCEDURE InitNormaliz 
   (   G  : Normaliz;
    Name  : ARRAY OF CHAR;
    Norm  : NormType;
    GType : GammaType;
    Beta  : LONGREAL;
    Maxt  : BDim;
       m  : MScal;
       k  : MDim
   );
\endcode
 \tab  Initializes the structure {\tt G} so that 
   {\tt Gamma[t]} contains the 
   approximation of $\gamma_t$ that corresponds to {\tt T},
   for $1\le t\le {\tt Maxt}$. 
   {\tt Cst[t]} will contain the value of $\bar\ell_t^2$ 
   in case of the Euclidean norm ({\tt Norm = L2Norm}), 
   and $\bar\ell_t$ otherwise,
   where $\bar\ell_t$ is defined in (\ref{eq:barellt}),
   $\tilde\gamma_t$ is determined by {\tt GType},  and
   $\beta$ is the value of {\tt Beta}.

   Normally, this procedure is called with $\beta = 1$.  
   Taking $\beta < 1$ inflates the figure of merit by
   $(1/\beta)^t$, thus weakening the requirements for large $t$ in a 
   worst-case figure of merit such as (\ref{eq:M}).

   The constant {\tt Maxt} cannot exceed 
   the maximum value of $t$ for which the approximation of $\gamma_t$
   has been pre-computed and is available (this depends on the type
   {\tt GType} of $\tilde\gamma_t$).
 \endtab
\code

PROCEDURE DeleteNormaliz (VAR G : Normaliz);
\endcode
 \tab  Deletes the structure {\tt G}.
 \endtab
\code

PROCEDURE WriteNormaliz (G : Normaliz);
\endcode
 \tab  Writes the information contained in {\tt G}.
 \endtab
\code

END NORM.
\endcode

